%% Section V-C through V-G draft: results prose.
%% Sources: Tables III-VI in draft_tables.tex; timing from
%% research/gtsam_poc/gtsam_poc_result_ps100_*.txt (6.0-9.0 ms/state, mean 7.8,
%% 42 s for the 87-minute line), x_abfullrelin_m4.log (68 ms/state),
%% x_abdense_m4.log (31 ms/state). Ground speed 68 m/s from the line exports.

\subsection{Representative line}\label{sec:repline}
Table~\ref{tab:1007} reports the longest line, 1007.06 (\SI{87}{\minute}), on
both cabin magnetometers, with the free-inertial drift at \SI{317.5}{\meter}.
Read causally, with no look-ahead at all, the proposed estimator reaches
\SI{42.7}{\meter} on Mag~4 and \SI{16.0}{\meter} on Mag~5, ahead of the
online-TL EKF (\num{46.7}/\SI{17.8}{\meter}) and of the NN-augmented filter
(\num{48.8}/\SI{18.3}{\meter}). The margin over the online-TL EKF is the purest
number in this paper: the two estimators share the error model, the
compensation model, the measurement model, and every noise setting, and differ
only in that the filter commits each state once while the graph relinearizes
and revises every state inside its lag. The margin over the EKF+TL+NN says
something further: the network exists to absorb what the linear TL basis
cannot, yet carrying the linear basis in a smoother beats carrying the linear
basis plus a network in a filter, on both sensors, without any learned
component. Allowed its \SI{300}{\second} lag, the committed estimate roughly
halves the causal error, to \num{24.2}/\SI{11.5}{\meter}. Fig.~\ref{fig:prior}
shows the error history.

\subsection{Breadth}\label{sec:breadth}
Table~\ref{tab:breadth} repeats the comparison on every counted line, spanning
two flights, two maps, two altitude regimes, and both cabin magnetometers, with
nothing re-tuned. The committed estimate is the most accurate entry on all
eight counted cases, between \SI{10.5}{} and \SI{28.4}{\meter}, and between
\num{1.5} and \num{5.2} times better than the strongest causal baseline on the
same case. Read causally, the proposed estimator is the best causal method on
six of the eight; it trails the EKF+TL+NN on 1007.02 Mag~5 by \SI{2.7}{\meter}
and on 1003.08 Mag~4 by \SI{2.3}{\meter}, and we report both rather than
claiming a sweep. Where the online-TL EKF fails outright, diverging on two
Mag~4 lines and leaving the map on a third, both readings of the proposed
estimator stay bounded; the mechanism is the one
Section~\ref{sec:sysmodel} anticipated, since the filter must explain the
uncompensated field with whatever state is cheapest at each instant, while the
smoother may revise its attribution for \SI{300}{\second} before committing.
The set-aside calibration line 1006.08 behaves as expected from its length and
content, and no conclusion is drawn from it.

\subsection{Lag design}\label{sec:lagdesign}
The lag is the estimator's single design variable, and Table~\ref{tab:lag}
sweeps it from \SI{30}{} to \SI{600}{\second} on the four counted lines. Two
regularities carry the design guidance. First, the committed error improves
with lag on every case, steeply up to about \SI{300}{\second} and weakly after
it: lengthening the lag from \SI{300}{} to \SI{600}{\second}, a
doubling of both latency and computation, buys at most \SI{2.6}{\meter} and on
one case slightly degrades, while the step from \SI{30}{} to \SI{300}{\second}
buys between \SI{3.6}{} and \SI{41.8}{\meter}. The knee reflects the physics: at
survey speed a \SI{300}{\second} lag corresponds to roughly \SI{20}{\kilo\meter}
of track, enough for the map gradient and the attitude-driven regressor to
decorrelate, after which further look-ahead adds little. We take
\SI{300}{\second} as the default. Second, the causal column is flat: across a
twentyfold change in lag the causal error moves by less than a metre on every
case. The lag decides only how long a state is revised before it is frozen, not
how good the newest state is, so the accuracy bought by look-ahead is entirely
smoothing gain, and a designer who cannot afford latency loses exactly the gap
between the two columns and nothing else.

\subsection{Ablation}\label{sec:ablation}
Table~\ref{tab:ablation} varies each element of the configuration on line
1007.06, one at a time. The estimator is insensitive at the metre level to all
of them. The compensation prior may move an order of magnitude in either
direction, or be replaced by a per-column prior that assigns each regressor
column \SI{100}{\nano\tesla} of measurement contribution, for changes under a
metre. Slowing the compensation walk a hundredfold costs about two metres,
consistent with a real installation state that drifts in flight; speeding it up
costs little. Removing the Huber kernel or widening its threshold costs one to
two metres, so robustness to map error helps but does not carry the result.
Inflating the assumed measurement noise fivefold, or relinearizing every
variable at every epoch, changes the estimate by less than half a metre; the
latter confirms that iSAM2's selective relinearization gives up nothing here.
The largest single effect is the process-noise floor: raising it from
$10^{-16}$ to $10^{-12}$ times the nominal position channel, i.e.\ from
$10^{-20}\mathbf{I}$ to $10^{-16}\mathbf{I}$, costs \SI{8.9}{\meter} on Mag~4,
which is why Section~\ref{sec:incremental} treats the floor as a conditioning
device to be kept minimal rather than a tuning knob. Attaching all ten raw
measurements to each state instead of one costs \SI{2.3}{\meter} on Mag~4 and
nothing on Mag~5, at four times the computation, supporting the
\SI{1}{\hertz} state cadence.

\subsection{Computational cost}\label{sec:cost}
All timings are single-core, on a laptop-class x86 processor under WSL2, with
the measurement factor evaluated through a Python callback; they are upper
bounds on a native implementation. At the reference configuration the
incremental update takes \SIrange{6.0}{9.0}{\milli\second} per state across the
ten line-magnetometer cases, mean \SI{7.8}{\milli\second}, so the
\SI{87}{\minute} line completes in \SI{42}{\second} and the margin to the
\SI{1}{\hertz} state cadence exceeds two orders of magnitude. The cost scales
with the number of retained states, so halving the lag halves it, and the two
expensive variants of the ablation bound the alternatives: relinearizing every
variable at every epoch costs \SI{68}{\milli\second} per state and attaching
all ten raw measurements \SI{31}{\milli\second}, both still real time at
\SI{1}{\hertz} but paying roughly ninefold and fourfold for accuracy changes
the ablation shows to be, respectively, below half a metre and a two-metre
degradation.
